% فصل اول: مقدمه و کلیات پژوهش
\chapter{مقدمه و کلیات پژوهش}
\label{ch:intro}

بسیاری از پدیده‌های بنیادی در علوم پایه و مهندسی، از جمله دینامیک سیالات، انتقال حرارت، مکانیک کوانتومی، انتشار امواج و فرایندهای بیوفیزیکی، در قالب معادلات دیفرانسیل جزئی\LTRfootnote{Partial Differential Equation (PDE)} مدل‌سازی می‌شوند. از این رو، ارائه روش‌های دقیق، کارآمد و محاسبات‌پذیر برای تحلیل و حل این معادلات همواره یکی از ارکان اصلی در حوزه محاسبات علمی بوده است. در دهه‌های اخیر، پیشرفت‌های شگرف در حوزه یادگیری عمیق\LTRfootnote{Deep Learning} تحولات بنیادینی را در پردازش تصویر، پردازش زبان طبیعی و علوم داده رقم زده است.

با این وجود، به‌کارگیری مستقیم مدل‌های یادگیری عمیق و صرفا داده‌محور (که به‌عنوان جعبه سیاه\LTRfootnote{Black Box} عمل نموده و صرفا همبستگی‌های آماری میان ورودی و خروجی را فرا می‌گیرند) در مسائل مهندسی و علوم فیزیکی با چالش‌های ساختاری مهمی روبه‌روست. در این گونه مسائل، مفهوم «داده‌های برچسب‌دار» به معنای اندازه‌گیری‌های تجربی از کمیت‌های میدان فیزیکی (نظیر سرعت سیال، توزیع فشار، دما یا غلظت گونه‌های شیمیایی) در نقاط مختلف مکانی و لحظات زمانی است. در بسیاری از سامانه‌های واقعی مهندسی — نظیر بررسی جریان سیال درون محفظه احتراق موتورهای جت، مخازن هیدروکربوری در اعماق زمین، رآکتورهای همجوشی هسته‌ای یا جریان خون درون رگ‌های زیستی — نصب حسگرهای متعدد و ثبت داده‌های پیوسته و متراکم، به دلایلی چون هزینه‌های گزاف آزمایشگاهی، شرایط محیطی خشن یا غیرقابل‌دسترس بودن هندسی، عملا ناممکن یا بسیار محدود است. بنابراین، مشاهدات موجود همواره به‌صورت تنک، پراکنده و همراه با خطای اندازه‌گیری هستند. در چنین شرایط کم‌داده‌ای، یک شبکه عصبی فاقد درک از قوانین حاکم، به‌شدت مستعد بیش‌برازش\LTRfootnote{Overfitting} بر روی نمونه‌های اندک بوده و در مناطقی که فاقد داده آموزشی است (نواحی برون‌یابی)، ممکن است پیش‌بینی‌هایی کاملا متناقض با بدیهیات فیزیک (مانند نقض قانون بقای جرم، پیش‌بینی مقادیر دمای منفی ناممکن، یا نقض شرایط مرزی سامانه) تولید نماید.

به‌منظور غلبه بر این محدودیت‌ها، ادغام دانش فیزیکی حاکم بر سیستم با ساختارهای یادگیری ماشین، رویکرد نوینی را تحت عنوان «شبکه‌های عصبی آگاه از فیزیک»\LTRfootnote{Physics-Informed Neural Network (PINN)} پدید آورده است. این چارچوب توسط رئیسی و همکارانش در سال ۲۰۱۹ معرفی شد \cite{raissi2019} و با استقبال گسترده جامعه علمی در مرز مشترک هوش مصنوعی و محاسبات علمی روبه‌رو گردید. در این رهیافت، معادله دیفرانسیل حاکم به‌همراه شرایط مرزی و اولیه مستقیما درون تابع هزینه مدل لحاظ می‌شود و شبکه به‌جای اتکای صرف به داده‌های اندازه‌گیری‌شده متراکم، قوانین فیزیک را در کل دامنه ارضا می‌نماید. این امر به شبکه امکان می‌دهد تا به عنوان یک حل‌گر عددی بدون شبکه محاسباتی عمل کند.

در این پایان‌نامه، چارچوب شبکه‌های عصبی آگاه از فیزیک برای حل مسئله مستقیم (محاسبه و بازسازی میدان پاسخ زمانی-مکانی) در معادله دیفرانسیل جزئی غیرخطی برگرز\LTRfootnote{Burgers Equation} مورد ارزیابی قرار می‌گیرد. معادله برگرز مدلی پایه‌ای در دینامیک سیالات و جریان ترافیک است که پدیده تشکیل «موج شوک»\LTRfootnote{Shock Wave} و ایجاد «جبهه‌های با شیب بسیار تند»\LTRfootnote{Steep Fronts / Steep Gradients} را توصیف می‌کند. در فیزیک، موج شوک به ناحیه‌ای بسیار باریک و فشرده گفته می‌شود که در آن مشخصات سیال (نظیر سرعت یا فشار) در فاصله مکانی بسیار کوتاهی دچار افت یا جهش ناگهانی و شدید می‌شوند (مانند پدیده شکستن امواج در ساحل یا متراکم شدن ناگهانی خودروها در ترافیک). از آن‌جا که توابع پیوسته و شبکه‌های عصبی استاندارد در تقریب‌زدن این‌گونه رفتارهای نوک‌تیز و شیب‌های تند با دشواری‌های زیادی روبه‌رو هستند، معادله برگرز به یک مسئله معیار و چالش‌برانگیز استاندارد برای ارزیابی دقت روش‌های عددی و یادگیری ماشین تبدیل شده است.

علاوه بر تحلیل نظری، پیاده‌سازی محاسباتی مدل زمان‌پیوسته با زبان برنامه‌نویسی پایتون و کتابخانه جکس\LTRfootnote{JAX} صورت گرفته است. جزئیات ساختار الگوریتمی، تبدیلات تابعی مشتق‌گیری خودکار، سازوکار کامپایل در زمان اجرا و استراتژی آموزش دومرحله‌ای به‌طور تفصیلی تدوین و اجرا شده‌اند. نتایج حاصل و کارایی مدل پیاده‌سازی‌شده در بازسازی جبهه شوک با مقاله مرجع معتبر رئیسی و همکاران (۲۰۱۹) مورد مقایسه و ارزیابی تطبیقی قرار گرفته است.

\section{طرح موضوع پژوهش}
\label{sec:subject}

موضوع محوری این پایان‌نامه «حل مسئله مستقیم معادله غیرخطی برگرز با استفاده از شبکه‌های عصبی آگاه از فیزیک» است. در مسئله مستقیم\LTRfootnote{Forward Problem}، ساختار و کلیه ضرایب فیزیکی معادله دیفرانسیل به‌همراه شرایط مرزی و اولیه به‌طور کامل مشخص بوده و هدف اصلی، تعیین توزیع مکانی-زمانی متغیر مجهول میدان پاسخ $u(t,x)$ با حداقل داده‌های اولیه و مرزی و بدون نیاز به تولید شبکه محاسباتی گسسته است.

رکن اساسی و موتور محرک محاسباتی در این چارچوب، بهره‌گیری از سازوکار «مشتق‌گیری خودکار»\LTRfootnote{Automatic Differentiation} است. در روش‌های محاسباتی سنتی، محاسبه مشتقات مکانی و زمانی معادلات دیفرانسیل مستلزم ایجاد شبکه محاسباتی و استفاده از تقریب‌های تفاضل محدود بود که با خطای برش و ناپایداری‌های ناشی از گام گسسته‌سازی همراه است. در مقابل، مشتق‌گیری خودکار تکنیکی بنیادین در علوم داده است که با بهره‌گیری از گراف محاسباتی برنامه و اعمال نظام‌مند قاعده زنجیره‌ای ریاضی، مشتقات جزئی خروجی شبکه نسبت به ورودی‌های فضا و زمان ($u_t, u_x, u_{xx}$) را با دقت تحلیلی و تا حد خطای ممیز شناور ماشین ارزیابی می‌کند. این قابلیت به مدل اجازه می‌دهد که بدون نیاز به شبکه‌بندی دامنه، باقیمانده فیزیکی معادله را در هر نقطه دلخواه محاسبه نموده و شبکه عصبی را وادار سازد که قوانین بقا و فیزیک حاکم را در سراسر دامنه پیوسته ارضا کند. مبانی نظری و سازوکار مشتق‌گیری خودکار به‌تفصیل در بخش \ref{sec:autodiff} تشریح خواهد شد.

\section{بیان مسئله}
\label{sec:problem}

روش‌های عددی کلاسیک نظیر روش تفاضل‌های محدود\LTRfootnote{Finite Difference Method} \cite{leveque2007}، روش اجزای محدود\LTRfootnote{Finite Element Method} \cite{zienkiewicz2005} و روش‌های طیفی\LTRfootnote{Spectral Methods} \cite{trefethen2000} طی چندین دهه ابزارهای استوار و تثبیت‌شده در حل معادلات دیفرانسیل جزئی بوده‌اند. با این حال، این روش‌ها به تولید شبکه محاسباتی وابستگی شدیدی دارند و در مسائل با ابعاد بالا، با پدیده نفرین ابعاد\LTRfootnote{Curse of Dimensionality} و رشد نمایی هزینه محاسباتی روبه‌رو می‌شوند. علاوه بر این، در مسائلی که شرایط هندسی نامنظم یا داده‌های مرزی پیچیده وجود دارد، اعمال روش‌های کلاسیک با دشواری‌های تکنیکی همراه است.

در نقطه مقابل، شبکه‌های عصبی استاندارد اگرچه از قابلیت بالایی در برازش توابع غیرخطی برخوردارند، اما در صورت عدم برخورداری از قیدهای ساختاری، به حجم بسیار زیادی از داده‌های آموزشی نیاز داشته و خروجی آن‌ها در مناطقی که فاقد داده آموزشی است، ممکن است دچار خطای برون‌یابی شدید و نقض آشکار قوانین فیزیکی گردد.

مسئله اصلی این پژوهش را می‌توان در این پرسش خلاصه کرد: چگونه می‌توان با تلفیق ساختار ریاضی معادله دیفرانسیل جزئی برگرز در فرایند آموزش شبکه‌های عصبی عمیق، یک ساختار محاسباتی کارآمد ایجاد کرد که بدون نیاز به شبکه‌بندی محاسباتی و صرفا با اتکا بر داده‌های مرزی-اولیه محدود و ارضای قید باقیمانده فیزیکی، میدان پیوسته پاسخ را با دقت بالا بازسازی کند و پدیده تشکیل موج شوک را با حداقل خطای عددی به دام اندازد؟ چارچوب شبکه‌های عصبی آگاه از فیزیک \cite{raissi2019} پاسخی ساختاریافته به این نیاز محاسباتی ارائه می‌دهد.

\section{اهداف پژوهش}
\label{sec:goals}

هدف کلی این پایان‌نامه، تحلیل مبانی نظری، تدوین الگوریتمی، پیاده‌سازی محاسباتی و ارزیابی تجربی چارچوب شبکه‌های عصبی آگاه از فیزیک در حل مسئله مستقیم معادله غیرخطی برگرز و مقایسه جامع نتایج با مراجع معتبر علمی است. اهداف تفصیلی پژوهش عبارت‌اند از:

\begin{enumerate}
\item بررسی و تدوین مبانی نظری معادلات دیفرانسیل جزئی، شبکه‌های عصبی پیش‌خور، روش‌های بهینه‌سازی و تکنیک‌های مشتق‌گیری خودکار،
\item تحلیل ساختار مدل زمان‌پیوسته شبکه‌های آگاه از فیزیک و فرمول‌بندی ریاضی تابع هزینه ترکیبی بر پایه باقیمانده معادله برگرز،
\item طراحی ساختار الگوریتمی و جریان اجرای محاسباتی مدل با بهره‌گیری از قابلیت‌های تبدیل تابعی و کامپایل در زمان اجرای کتابخانه جکس،
\item پیاده‌سازی محاسباتی مدل زمان‌پیوسته برای حل مستقیم معادله برگرز با استفاده از کتابخانه جکس و الگوریتم بهینه‌سازی ترکیبی آدام و \lr{L-BFGS}،
\item ارزیابی رفتار همگرایی توابع هزینه، تحلیل توزیع مکانی-زمانی خطا و سنجش تطبیقی نتایج با مقاله معتبر مرجع رئیسی و همکاران (۲۰۱۹).
\end{enumerate}

\section{پرسش‌های پژوهش}
\label{sec:questions}

این پژوهش در صدد پاسخ‌گویی به پرسش‌های علمی زیر است:

\begin{enumerate}
\item آیا یک شبکه عصبی با ساختار استاندارد قادر است تنها با اتکا بر داده‌های مرزی-اولیه محدود و باقیمانده معادله، پاسخ معادله دیفرانسیل جزئی غیرخطی برگرز را در کل دامنه زمانی-مکانی بدون تولید شبکه محاسباتی تقریب بزند؟
\item فرایند مشتق‌گیری خودکار چگونه باقیمانده فیزیکی معادله برگرز را در ساختار تابع هزینه مدل تلفیق می‌نماید؟
\item ساختار الگوریتمی و پیاده‌سازی تابعی در کتابخانه جکس چه تاثیری در سرعت و کارایی محاسباتی آموزش مدل ایفا می‌کند؟
\item الگوریتم بهینه‌سازی ترکیبی دو مرحله‌ای (آدام و \lr{L-BFGS}) چه نقشی در سرعت و دقت همگرایی مدل ایفا می‌کند؟
\item عملکرد و دقت مدل پیاده‌سازی‌شده در مقایسه با نتایج گزارش‌شده در مقاله معتبر مرجع \cite{raissi2019} چگونه ارزیابی می‌شود؟
\end{enumerate}

\section{روش‌شناسی پژوهش}
\label{sec:methodology}

روش‌شناسی به کار رفته در این پایان‌نامه مبتنی بر تلفیق مطالعات تحلیلی-نظری و پیاده‌سازی‌های محاسباتی-تجربی است:

در گام نخست، پیشینه پژوهش و منابع معتبر علمی در حوزه‌های معادلات دیفرانسیل جزئی \cite{evans2010}، روش‌های عددی کلاسیک \cite{leveque2007,zienkiewicz2005}، یادگیری عمیق \cite{goodfellow2016} و مشتق‌گیری خودکار \cite{baydin2018} مورد بررسی و تحلیل جامع قرار گرفت و مبانی نظری شبکه باقیمانده و توابع هزینه تلفیقی استخراج شد.

در گام دوم، معادله برگرز یک‌بعدی با شرایط اولیه سینوسی و شرایط مرزی همگن به‌دلیل رفتار موج شوک غیرخطی به‌عنوان مسئله آزمون انتخاب شد.

در گام سوم، مدل زمان‌پیوسته شبکه عصبی آگاه از فیزیک با استفاده از معماری پرسپترون چندلایه، تبدیلات تابعی مشتق‌گیری و کامپایل بهینه‌ساز در زبان برنامه‌نویسی پایتون با کتابخانه محاسباتی جکس پیاده‌سازی شد.

در گام چهارم، آزمایش‌های عددی حل مسئله مستقیم با بهینه‌سازی ترکیبی دومرحله‌ای (آدام و سپس \lr{L-BFGS}) اجرا گردید و تاریخچه همگرایی مولفه‌های هزینه، بازسازی جبهه شوک و معیارهای دقتی در مقایسه با مقاله مرجع رئیسی و همکاران (۲۰۱۹) مورد ارزیابی قرار گرفت.

\section{ساختار پایان‌نامه}
\label{sec:structure}

فصول این پایان‌نامه به‌صورت زیر تدوین شده است:
\begin{itemize}
\item \textbf{فصل دوم (مبانی نظری و پیشینه پژوهش):} مروری بر مفاهیم پایه‌ای معادلات دیفرانسیل جزئی، روش‌های عددی کلاسیک، ساختار شبکه‌های عصبی عمیق، روش‌های آموزش و بهینه‌سازی و بررسی تاریخی پژوهش‌های پیشین در یادگیری ماشین فیزیک‌پایه.
\item \textbf{فصل سوم (شبکه‌های عصبی آگاه از فیزیک):} تشریح جامع مبانی مشتق‌گیری خودکار و حساب اعداد دوگانه، تبیین دقیق نحوه به‌کارگیری مشتق‌گیری خودکار در استخراج باقیمانده معادله برگرز، فرمول‌بندی مدل‌های زمان‌پیوسته و زمان‌گسسته و نکات عملی در آموزش شبکه.
\item \textbf{فصل چهارم (پیاده‌سازی و ارزیابی نتایج):} ارائه جزئیات ساختار الگوریتمی و پیاده‌سازی با کتابخانه جکس، تحلیل روند آموزش و همگرایی مدل، ارزیابی بازسازی میدان پاسخ و مقایسه جامع یافته‌ها با مقاله معتبر مرجع.
\item \textbf{فصل پنجم (نتیجه‌گیری و پیشنهادها):} جمع‌بندی دستاوردهای پژوهش، تبیین محدودیت‌های موجود و ارائه چشم‌اندازها و پیشنهادهایی برای پژوهش‌های آتی.
\end{itemize}
